\documentclass[preprint,12pt]{elsarticle}

\usepackage{amsmath,amssymb}
\usepackage{booktabs}
\usepackage{siunitx}
\usepackage{graphicx}
\usepackage{hyperref}
\usepackage{float}
\usepackage{xurl}
\usepackage{microtype}
\setlength{\emergencystretch}{2em}
\graphicspath{{figures/}}

\journal{Transportation Research Part E: Logistics and Transportation Review}

\begin{document}
\begin{frontmatter}

\title{Initial-Action Envelope for Pilot Man-Overboard Response: Stern-Clearance Time and Heading Change}

\author[inst1,inst4]{Gokhan Camliyurt\fnref{fn2}}
\ead{au810644@uni.au.dk}
\author[inst5]{Juliet Wamugi\fnref{fn1}}
\ead{jwangui@jkuat.ac.ke}
\author[inst1]{Zaili Yang\corref{cor1}}
\ead{Z.Yang@ljmu.ac.uk}

\cortext[cor1]{Corresponding author}
\fntext[fn1]{ORCID: 0009-0002-5174-7780}
\fntext[fn2]{ORCID: 0009-0001-7762-0932}

\address[inst1]{Liverpool Logistics, Offshore and Marine Research Institute, Liverpool John Moores University, Liverpool, United Kingdom}
\address[inst4]{Department of Navigation Sciences, National Korea Maritime and Ocean University, Busan, South Korea}
\address[inst5]{Jomo Kenyatta University of Agriculture and Technology, Nairobi, Kenya}

\begin{abstract}
Pilot man-overboard events during transfer operations require immediate bridge action. In this first phase, the objective is rapid separation of the casualty from stern-propeller danger, before any full recovery-turn pattern is completed. Existing guidance remains largely qualitative and does not provide event-wise hold-time guidance for initial hard-over action. This study develops a practical stern-clearance model and applies it to 5,302 eligible scenarios reconstructed from operational event data. The model combines low-order yaw response, bridge and rudder timing, stern geometry, loading effects, and environmental drift. Results show three clear outcome groups: no hazard entry, clearance after entry, and no clearance within the simulation window. For clearance-after-entry cases, the central tendency is about one minute, while upper-tail cases require substantially longer action durations. The study provides a reproducible baseline for initial-action decision support in pilot MOB response, with detailed limitations and data-availability constraints discussed in the manuscript.
\end{abstract}

\begin{keyword}
Man overboard \sep Crew-witnessed MOB \sep Stern-clearance time \sep Ship maneuvering \sep Decision support
\end{keyword}

\end{frontmatter}

\section{Introduction}
Pilot transfer is a high-risk operation in constrained approaches. In a pilot MOB case, the first tactical objective is immediate stern-propeller separation, not full recovery-turn completion. However, bridge practice still relies on qualitative commands (e.g., hard-over toward casualty side), with limited quantitative guidance on required hard-rudder duration.

This study formulates immediate protection as a time-to-clearance problem and estimates operational thresholds for the first pilot MOB action phase after detection. Here, crew-witnessed MOB denotes a witnessed event that demands immediate bridge action. The objective is not to reconstruct the full emergency turn. Instead, the paper isolates initial protection quantities that can be used as trigger thresholds. The framework explicitly represents bridge-response delay, rudder-travel time, vessel maneuvering response, and stern-hazard geometry, while retaining environmental forcing when compatible metocean magnitudes are available.

\subsection*{Contributions}
This paper contributes a physically interpretable first-phase protection metric, a paired trigger-variable framework for immediate action, a data-feasible model with explicit fallback handling for missing maneuver or metocean inputs, and percentile guidance stratified by ship type, LOA bin, and fall side.

\section{Literature review}
\subsection{MOB doctrine and first-phase decision gap}
Classical MOB procedures emphasize recovery geometry, track reversal, and retrieval completion. In pilot-transfer cases, bridge action is needed before any full recovery pattern can be established. The near-term objective is rapid separation between the casualty and the stern-propeller danger region. The regulatory baseline (SOLAS and IAMSAR) requires immediate MOB action, but does not provide event-wise quantitative hold-time thresholds for initial hard-over protection \cite{ref_solas,ref_iamsar}. Recent pilot-transfer accident evidence confirms that this phase is highly time-critical under coupled environmental, boarding, and organizational pressures \cite{ref_pilot_transfer_accidents}. Recent pilot-focused studies also show that pilot-transfer accidents are driven by interacting human, environmental, and vessel factors, and that operational risk modelling for pilot accidents is feasible with both tree-based and machine-learning approaches \cite{ref_camliyurt_pilot_fta2022,ref_camliyurt_pilot_ml2023}. Existing doctrine gives practical commands, but rarely quantifies the duration of initial hard-rudder action or the heading change needed for first-phase protection.

\subsection{Low-order maneuvering models in operational settings}
Low-order maneuvering models remain operationally useful when vessel-specific hydrodynamic derivatives, trial curves, or high-fidelity simulation inputs are unavailable. Recent reviews show that maritime decision support now spans physical, reduced-order, data-driven, and hybrid approaches, with growing emphasis on short-horizon safety prediction \cite{ref_ship_motion_review,ref_ship_motion_review2}. For event-wise operational analysis, first-order yaw-response formulations are attractive because they maintain physical interpretability while enabling large-sample scenario evaluation; the standard Nomoto-type interpretation is well established in marine craft motion-control literature \cite{ref_fossen}. Their limitations are established, particularly for nonlinear rudder effects, propulsion transients, and shallow-water interactions. Even so, for early protection questions posed as short-time envelope estimation, reduced-order dynamics remain an appropriate and reproducible baseline.

\subsection{AIS-based safety analytics}
AIS-based safety analysis is now central to encounter-risk assessment, traffic-exposure quantification, port-approach analysis, and transfer monitoring. Recent reviews show extensive use of AIS trajectories for anomaly detection, behavior interpretation, and decision-support feature extraction \cite{ref_ais_review}. In pilotage research, AIS data are valuable for reconstructing the operational context of low-frequency, high-consequence events. However, most AIS studies remain descriptive at the traffic-pattern level. Fewer studies derive explicit first-response maneuver thresholds, and translation to stern-hazard time-to-clearance metrics remains limited.

\subsection{Gap positioning}
Taken together, the maritime literature indicates three unresolved needs: no standardized quantitative threshold for immediate stern-clearance protection in pilot-transfer MOB cases, limited explicit modeling of the bridge-response chain and rudder actuation delay that govern the earliest event seconds, and lack of percentile guidance stratified by ship type, size, and transfer side.

In parallel, maritime safety-risk literature has developed mature Bayesian-network frameworks for accident and inspection risk modelling, including influential contributions by Zaili Yang and colleagues \cite{ref_yang_bn_maritime2014,ref_yang_psc_bn2018,ref_yang_maritime_accident_bn2020}. These studies provide strong methodological grounding for probabilistic risk structuring, but they do not directly resolve the immediate pilot MOB stern-clearance trigger problem addressed here.

Cross-disciplinary transport research reinforces this framing. In aviation, recent findings show that pilot response to cockpit alerts is variable, context-dependent, and tightly linked to warning-system design \cite{ref_aviation_alerts,ref_aviation_hud}. Aviation control studies further formalize safety as envelope protection, i.e., maintaining operation within a dynamic safe region rather than reacting only after severe state deviation \cite{ref_aviation_envelope}. In land transportation, emergency-braking and collision-avoidance research similarly uses short-horizon trigger variables such as time-to-collision, detection timing, and braking onset \cite{ref_land_aeb,ref_land_detection}. These domains do not transfer directly to maritime MOB mechanics, but they provide convergent methodological support for modeling the first action as an envelope-and-trigger problem prior to full maneuver optimization.

The present study addresses these gaps by isolating the first emergency action phase and expressing it through paired trigger variables: clearance time and heading at clearance.

\section{Methodology}
Figure~\ref{fig:research_flow} outlines the five-step workflow adopted in this study, from problem framing and variable design through event-wise scenario construction, batch simulation, and percentile guidance derivation. The subsections below detail each step.

\begin{figure}[H]
    \centering
    \includegraphics[width=0.95\textwidth]{Figure_1_research_flow.png}
    \caption{Methodology workflow.}
    \label{fig:research_flow}
\end{figure}

\subsection{Nomenclature}
Table~\ref{tab:nomenclature} defines the core symbols used throughout the paper. All timing variables carry units of seconds; angular variables are reported in degrees in all results tables and figures.
\begin{table}[H]
\centering
\caption{Core symbols and units.}
\label{tab:nomenclature}
\begin{tabular}{lll}
\toprule
Symbol & Meaning & Unit \\
\midrule
$\tau_{\mathrm{clear}}$ & Stern-clearance time & s \\
$\psi$ & Heading response & rad/deg \\
$r$ & Yaw rate & rad/s \\
$K$ & Steering gain (rudder effectiveness) & 1/s \\
$T_r$ & Yaw time constant (response lag) & s \\
$\delta$ & Commanded rudder angle (hard-over) & deg \\
$\delta_{\mathrm{eff}}$ & Effective rudder angle after ramp & deg \\
$t_{\mathrm{response}}$ & Total bridge/helm delay & s \\
$t_{\mathrm{rudder}}$ & Rudder travel time (midship to hard-over) & s \\
$l_s$ & Longitudinal distance from pivot to propeller & m \\
$x_s$ & Longitudinal distance from midship to pivot & m \\
$R_h$ & Effective hazard radius around propeller & m \\
$D_p$ & Propeller diameter & m \\
$\mathbf{r}_c(t)$ & Casualty position vector in ship-fixed frame & m \\
$\mathbf{r}_p(t)$ & Stern/propeller centre position vector & m \\
$u_c,\,v_c$ & Casualty velocity components (surge, sway) & m/s \\
$\lambda_d$ & Draft ratio ($\text{draft}/\text{LOA}$) & -- \\
\bottomrule
\end{tabular}
\end{table}

\subsection{Clearance definition}
The primary output variable is the stern-clearance time, defined as the first instant at which the casualty exits the propeller-hazard zone after having entered it:
\begin{equation}
\tau_{\mathrm{clear}} = \inf\left\{t>0:\ \|\mathbf{r}_c(t)-\mathbf{r}_p(t)\|\ge R_h\right\},
\end{equation}
where $\mathbf{r}_c(t)$ is the position of the casualty in the ship-fixed frame, $\mathbf{r}_p(t)$ is the position of the propeller centre as the stern swings, and $R_h$ is the effective hazard radius around the propeller. The condition $\|\mathbf{r}_c-\mathbf{r}_p\|\ge R_h$ is evaluated only after a prior entry ($\|\mathbf{r}_c-\mathbf{r}_p\|<R_h$), making $\tau_{\mathrm{clear}}$ the first outward crossing. This inf-type formulation is analogous to time-to-collision and envelope-exit metrics used in adjacent safety domains \cite{ref_land_aeb,ref_aviation_envelope}.

\subsection{Response-chain timing}
The ship-response clock is decomposed into a bridge-response block and a rudder-actuation block. The bridge-response delay is written as
\begin{equation}
t_{\mathrm{bridge}} = t_{\mathrm{detect/report}} + t_{\mathrm{master\ decision}} + t_{\mathrm{order}} + t_{\mathrm{helmsman}},
\end{equation}
and the total ship-response lag becomes
\begin{equation}
t_{\mathrm{ship\ response}} = t_{\mathrm{bridge}} + t_{\mathrm{rudder}}.
\end{equation}
This decomposition follows the same perception--decision--actuation logic used in adjacent transportation response studies, including cockpit alert handling and road-driver emergency response \cite{ref_aviation_alerts,ref_land_detection,ref_land_aeb}. In the present baseline, $t_{\mathrm{bridge}}=4.5$~s is adopted as a fixed high-readiness assumption, representing a best-case scenario where all bridge personnel are alert and the command chain executes without hesitation; this value is consistent with short-notice emergency response benchmarks in maritime bridge operations \cite{ref_solas,ref_iamsar}. The rudder travel time $t_{\mathrm{rudder}}\approx 15.08$~s is derived from the IMO steering-gear benchmark proxy (midship to $35^{\circ}$ at the prescribed maximum rate of $2.3^{\circ}$/s under SOLAS Chapter II-1) \cite{ref_solas}.

The casualty-exposure clock is not prescribed as a separate fixed delay. Instead, time to hazard entry emerges from the relative motion between casualty and moving stern hazard. For reference, an entry time can be written as
\begin{equation}
\tau_{\mathrm{entry}} = \inf\left\{t>0:\ \|\mathbf{r}_c(t)-\mathbf{r}_p(t)\| < R_h\right\},
\end{equation}
which makes clear that pilot fall-to-hazard timing is determined by initial position, ship speed, environmental drift, and stern swing, rather than by a separate scalar assumption.

\subsection{Yaw and rudder dynamics}
Yaw response follows the first-order Nomoto model \cite{ref_fossen}:
\begin{align}
\dot{\psi}(t) &= r(t),\\
\dot{r}(t) &= \frac{K\,\delta_{\mathrm{eff}}(t)-r(t)}{T_r}.
\end{align}
Here $K$ is the steering gain, reflecting how effectively a unit rudder angle translates into yaw-rate response (higher $K$ means a more agile vessel), and $T_r$ is the yaw time constant, reflecting the lag before the vessel attains the steady-state yaw rate (higher $T_r$ means a sluggish response). The model is the simplest physically interpretable representation of ship course-changing dynamics and is appropriate for the short-horizon enveloping problem studied here; nonlinear rudder saturation and shallow-water corrections are outside the scope of this baseline.

The effective rudder angle $\delta_{\mathrm{eff}}(t)$ is modelled as a three-phase ramp to represent the mechanical delay and finite travel rate of the helm actuation system:
\begin{equation}
\delta_{\mathrm{eff}}(t)=
\begin{cases}
0, & t\le t_{\mathrm{response}},\\
\delta\dfrac{t-t_{\mathrm{response}}}{t_{\mathrm{rudder}}}, & t_{\mathrm{response}}<t<t_{\mathrm{response}}+t_{\mathrm{rudder}},\\
\delta, & t\ge t_{\mathrm{response}}+t_{\mathrm{rudder}}.
\end{cases}
\end{equation}
The first branch (zero rudder) spans the bridge-response delay; the second branch models a linear ramp from midship to hard-over, approximating constant-rate helm travel consistent with the IMO steering-gear requirement \cite{ref_solas}; the third branch holds the commanded hard-over angle constant once full deflection is reached.

\subsection{Geometry and drift}
As the vessel turns through heading angle $\psi(t)$, the propeller centre traces an arc in the ship-fixed frame. Its coordinates relative to the initial midship origin are:
\begin{align}
 x_p(t) &= -x_s + l_s\left(1-\cos\psi(t)\right),\\
 y_p(t) &= -l_s\sin\psi(t),
\end{align}
where $l_s$ is the longitudinal distance from the vessel's pivot point to the propeller centre, and $x_s$ is the longitudinal distance from midship to the pivot point. The cosine term captures the forward retraction of the stern as the ship turns, and the sine term captures the lateral sweep of the propeller. Together these define the moving hazard centre $\mathbf{r}_p(t)=(x_p(t),\,y_p(t))$.

The casualty is modelled as a point drifting at constant velocity in the ship-fixed frame:
\begin{equation}
\mathbf{r}_c(t)=\mathbf{r}_c(0)+\begin{bmatrix}u_c\\v_c\end{bmatrix}t,
\end{equation}
where $u_c$ and $v_c$ are the net surge and sway velocity components of the casualty relative to the ship, combining the ship's forward speed, environmental drift (projected wind, wave, and current forcing), and the initial transfer-board momentum. The constant-drift assumption is appropriate for the short time horizons examined ($t_{\max}=180$~s) and is consistent with short-horizon drift models used in search-and-rescue trajectory prediction. The ship-fixed frame is defined as $x$ forward and $y$ starboard, with the casualty's initial coordinates referenced to the transfer point.

\subsection{Loading-state proxy}
When sea-trial derivatives are unavailable, the loading condition is approximated by the draft ratio $\lambda_d=\mathrm{draft}/\mathrm{LOA}$. Loaded vessels are generally less manoeuvrable (lower $K$, higher $T_r$) than vessels in ballast, consistent with physical expectations and documented in standard maneuvering references \cite{ref_fossen}. Three loading bands are applied with scalar multipliers:
\begin{itemize}
    \item ballast-like ($\lambda_d$ low): $K\times1.10$, $T_r\times0.85$ --- more responsive;
    \item mid-load: $K\times1.00$, $T_r\times1.00$ --- nominal;
    \item loaded-like ($\lambda_d$ high): $K\times0.90$, $T_r\times1.15$ --- less responsive.
\end{itemize}

\subsection{Data inputs and baseline assumptions}
The executable dataset combines four input groups: an AIS-derived event table (event identifiers, vessel particulars, speed, heading, and transfer side), a maneuver-plan table (rudder command records), metocean-attached fields (wind, wave, and current direction/magnitude when available), and a scenario-builder output table used to complete missing fields and apply documented fallback rules. Table~\ref{tab:data_mapping} maps these sources to their roles in the baseline run.

\begin{table}[H]
\centering
\caption{Executable inputs and variable mapping for the baseline run.}
\label{tab:data_mapping}
\begin{tabular}{p{0.25\textwidth}p{0.30\textwidth}p{0.37\textwidth}}
\toprule
Source & Key fields & Model role \\
\midrule
Accessible event table from current project source folders & event id, LOA, beam, draft, speed, heading, side & Base event state for the executable baseline \\
Accessible maneuver-plan table from current project source folders & emergency rudder command & Used when temporally matched; otherwise replaced by documented hard-over fallback \\
Metocean-attached fields in executable event table & wind/wave/current direction and magnitude & Optional drift forcing when magnitudes are present in the executed rows \\
Scenario builder & delay, rudder ramp, hazard radius, loading proxy & Event completion, eligibility filtering, and fallback bookkeeping \\
\bottomrule
\end{tabular}
\end{table}

Table~\ref{tab:baseline_assumptions} documents the scalar assumptions applied when event-level inputs are absent or unmatched. Key fixed values are the bridge-response delay ($t_{\mathrm{bridge}}=4.5$~s, high-readiness assumption) and the rudder travel time ($t_{\mathrm{rudder}}\approx15.08$~s from the IMO steering-gear benchmark proxy). Two additional delay cases ($t_{\mathrm{bridge}}+10$~s and $t_{\mathrm{bridge}}+30$~s) are examined in the one-at-a-time sensitivity analysis to bound the sensitivity of results to this fixed assumption.

\begin{table}[H]
\centering
\caption{Baseline assumptions used in the current implementation.}
\label{tab:baseline_assumptions}
\begin{tabular}{p{0.33\textwidth}p{0.50\textwidth}p{0.12\textwidth}}
\toprule
Item & Assumption & Value/range \\
\midrule
Initial longitudinal casualty position & Midship transfer location & $x_{c0}=0$ \\
Fall-side convention & Starboard $\Rightarrow y_{c0}>0$, port $\Rightarrow y_{c0}<0$ & by event \\
Rudder cap & Commanded hard-over upper bound & $35^{\circ}$ \\
Rudder source in current run & No matched maneuver-plan record; fallback baseline & $35^{\circ}$ fixed \\
Bridge/helm delay & Detection/report + decision + order + helm action & $\sim4.5$ s \\
Bridge/helm delay variability & Baseline is fixed in current run; variability assessed by sensitivity runs & +10 s, +30 s cases \\
Rudder travel time & IMO steering-gear benchmark proxy (0$^{\circ}$ to 35$^{\circ}$) & $\sim15.08$ s \\
Hazard radius & Effective envelope scaling & $R_h\approx2.5D_p$ \\
Metocean forcing in current run & Wind and wave bridged for most eligible events; current unavailable & partially active \\
\bottomrule
\end{tabular}
\end{table}

In this baseline, $R_h\approx2.5D_p$ is treated as a conservative engineering envelope parameter (not a regulatory threshold), and its influence is evaluated in the one-at-a-time sensitivity analysis.

\subsection{Model coverage checklist}
Figure~\ref{fig:algorithm_pipeline} presents the full decision-gated algorithm flow from MOB event initialisation through eligibility filtering, response-chain timing, rudder actuation, yaw dynamics, and outward-crossing evaluation. Left branches represent BAN or unresolved outcomes; right branches represent ALLOW outcomes.

\begin{figure}[H]
    \centering
    \includegraphics[width=0.85\textwidth]{Figure_2_algorithm_pipeline.png}
    \caption{Decision-gated algorithm flow for stern-clearance computation.}
    \label{fig:algorithm_pipeline}
\end{figure}

Figure~\ref{fig:model_architecture} decomposes the same logic into four interacting model blocks---maneuvering, geometry, environment, and hazard/crossing---with arrows indicating information flow toward the output variables $\tau_{\mathrm{clear}}$ and $\psi_{\mathrm{clear}}$.

\begin{figure}[H]
    \centering
    \includegraphics[width=0.92\textwidth]{Figure_3_model_architecture.png}
    \caption{Model architecture and information flow between maneuvering, geometry, environment, and hazard blocks.}
    \label{fig:model_architecture}
\end{figure}

Table~\ref{tab:model_coverage_checklist} provides a row-by-row traceability mapping between the algorithm pipeline and the implemented calculation blocks, indicating inclusion status and flagging items planned for future extension.

\begin{table}[H]
\centering
\caption{Figure-to-algorithm coverage checklist for MOB stern-clearance computation}
\label{tab:model_coverage_checklist}
\begin{tabular}{p{0.43\textwidth}p{0.13\textwidth}p{0.13\textwidth}p{0.22\textwidth}}
    \toprule
Calculation block & In Fig. 2 & In code & Status \\
\midrule
MOB start and initial position ($t=0$, $x_{c0}=0$, side by $y_{c0}$ sign) & Yes & Yes & Included \\
Response chain ($t_{\mathrm{detect}},t_{\mathrm{report}},t_{\mathrm{master}},t_{\mathrm{order}},t_{\mathrm{helm}}$) and $t_{\mathrm{response}}$ & Yes & Yes & Included \\
Rudder command and actuation ($\delta$, $t_{\mathrm{rudder}}$, $\delta_{\mathrm{eff}}(t)$) & Yes & Yes & Included, fallback active \\
Yaw maneuvering equations ($\dot{\psi}=r$, $\dot{r}=(K\delta_{\mathrm{eff}}-r)/T_r$) & Yes & Yes & Included \\
Loading proxy ($\lambda_d=\mathrm{draft}/\mathrm{LOA}$) and $(K,T_r)$ multipliers & Yes & Yes & Included \\
Wind and current directional projection to body frame & Yes & Yes & Wind active; current unavailable in current run \\
Wave influence in drift term (height and direction proxy) & Yes & Yes & Active for most eligible events \\
Casualty and propeller-hazard kinematics ($\mathbf{r}_c(t)$, $\mathbf{r}_p(t)$, $R_h$) & Yes & Yes & Included \\
Hazard entry and first outward-crossing logic for $\tau_{\mathrm{clear}}$ & Yes & Yes & Included \\
No-entry and no-crossing branch outcomes ($\tau_{\mathrm{clear}}=0$ or unresolved) & Yes & Yes & Included \\
Percentile aggregation ($P10,P50,P90,P95$) for eligible events & Yes & Yes & Included \\
Propeller RPM/thrust-state coupling at MOB instant & No & No & Missing (planned) \\
Wave period/spectral descriptor (beyond height-direction proxy) & No & No & Missing (planned) \\
Uncertainty propagation (e.g., Monte Carlo) on delays and drift coefficients & No & No & Missing (planned) \\
\bottomrule
\end{tabular}
\end{table}

\section{Results and discussion}
The executable analysis retains operationally eligible cases in the category ``allow with emergency plan. This yields 5{,}302 scenarios, of which 5{,}270 return solved outcomes. In the current run, no maneuver-plan timestamp matches eligible events; therefore, all events apply the documented $35^{\circ}$ hard-over fallback. Baseline bridge-response delay is fixed at 4.5 s for all events (high-readiness assumption), while rudder travel uses an IMO steering-gear benchmark proxy of about 15.08 s from midship to $35^{\circ}$. Environmental forcing is partially active: wind and wave fields are attached for most eligible events, while current remains unavailable in the executable baseline.

\subsection{Outcome classes and event disposition}
The solved set is divided into explicit physical outcome classes. Of 5{,}302 events, 2{,}010 (37.91\%) are \emph{maneuver effective -- no hazard entry} (never-entered), 3{,}260 (61.49\%) are \emph{clearance after entry} (cleared-after-entry), and 32 (0.60\%) are \emph{entered but not cleared within simulation horizon} (entered-no-clear). The no-entry class is a successful protection outcome and is not treated as missing data.
The 32 no-clearance events form a coherent high-risk cluster (near-dead-slow SOG, low steering gain, and predominantly port-side falls), indicating physically constrained maneuver effectiveness rather than numerical failure. Event counts and proportions for all three classes are reported in Table~\ref{tab:outcome_classes}.

\begin{table}[H]
\centering
\caption{Outcome-class summary for eligible scenarios.}
\label{tab:outcome_classes}
\begin{tabular}{lrr}
	\toprule
Outcome class & N & \% of total \\
\midrule
Maneuver effective -- no hazard entry (never-entered) & 2{,}010 & 37.91 \\
Clearance after entry (cleared-after-entry) & 3{,}260 & 61.49 \\
Entered but no clearance by $t_{\max}$ (entered-no-clear) & 32 & 0.60 \\
\bottomrule
\end{tabular}
\end{table}

\subsection{Overall pilot MOB envelope}
For clearance-after-entry events, the overall $\tau_{\mathrm{clear}}$ envelope is $P10=16.4$ s, $P50=63.1$ s, $P90=133.2$ s, and $P95=144.8$ s (Table~\ref{tab:main_percentiles}; see Figure~\ref{fig:tau_distribution} for the full distribution with percentile markers overlaid). The distribution therefore indicates a strongly spread operational range: some cases clear the stern-hazard zone within a relatively short interval, whereas others require substantially longer action duration before outward crossing occurs. This spread is consistent with the intended interpretation of the study variable. The quantity is not a single deterministic threshold valid for all transfer events; rather, it is an operational envelope whose upper percentiles represent more conservative initial-action guidance.

\begin{table}[H]
\centering
\caption{Operational percentile summary (nonzero $\tau_{\mathrm{clear}}$)}
\label{tab:main_percentiles}
\begin{tabular}{lrrrr}
    \toprule
Scope & P10 (s) & P50 (s) & P90 (s) & P95 (s) \\
\midrule
Eligible nonzero events ($n=3{,}260$) & 16.4 & 63.1 & 133.2 & 144.8 \\
Port-side nonzero events ($n=2{,}590$) & 19.0 & 83.0 & 137.22 & 147.51 \\
Starboard-side nonzero events ($n=670$) & 11.8 & 28.0 & 49.8 & 56.51 \\
\bottomrule
\end{tabular}
\end{table}

\begin{figure}[H]
    \centering
    \includegraphics[width=0.92\textwidth]{Figure_4_tau_distribution.png}
    \caption{Distribution of nonzero stern-clearance times.}
    \label{fig:tau_distribution}
\end{figure}

For the paired heading trigger variable in clearance-after-entry events, $\psi_{\mathrm{clear}}$ percentiles are $P10=0.3648^{\circ}$, $P50=6.9835^{\circ}$, $P90=16.7785^{\circ}$, and $P95=18.7623^{\circ}$ (Table~\ref{tab:psi_percentiles}).
These heading-change magnitudes are comparatively small, implying that heading change alone is a weaker bridge-level proxy for immediate protection than time-based guidance; in practice, $\tau_{\mathrm{clear}}$ is the more robust operational trigger.

Side-wise nonzero envelopes are strongly asymmetric: for port-side cases ($n=2{,}590$), the percentiles are $P10=19.0$ s, $P50=83.0$ s, $P90=137.22$ s, and $P95=147.51$ s; for starboard-side cases ($n=670$), they are $P10=11.8$ s, $P50=28.0$ s, $P90=49.8$ s, and $P95=56.51$ s.

\begin{table}[H]
\centering
\caption{Heading-at-clearance summary for clearance-after-entry events}
\label{tab:psi_percentiles}
\begin{tabular}{lrrrr}
    \toprule
Scope & P10 (deg) & P50 (deg) & P90 (deg) & P95 (deg) \\
\midrule
Eligible clearance-after-entry events ($n=3{,}260$) & 0.3648 & 6.9835 & 16.7785 & 18.7623 \\
Port-side clearance-after-entry events ($n=2{,}590$) & 0.6027 & 9.0135 & 17.401 & 19.1526 \\
Starboard-side clearance-after-entry events ($n=670$) & 0.1858 & 1.4195 & 4.1856 & 4.9994 \\
\bottomrule
\end{tabular}
\end{table}

\subsection{Interpretation of side asymmetry}
The side-stratified results show a pronounced asymmetry between port and starboard cases in the current baseline (see also Figure~\ref{fig:stratified_heatmap} for stratification by ship type and size). Port-side cases exhibit much larger median and upper-percentile clearance times than starboard cases. Under the present sign convention, positive hard-over produces positive yaw and stern swing toward negative-$y$ (port in the body frame), so port-side casualties are more often in the stern-swing path while starboard-side casualties are more often protected by outward separation. This mechanism is consistent with the observed disposition split (port: 1{,}378 never-entered, 2{,}590 cleared-after-entry, 31 no-clear; starboard: 632 never-entered, 670 cleared-after-entry, 1 no-clear).

\begin{figure}[H]
    \centering
    \includegraphics[width=0.98\textwidth]{Figure_5_stratified_heatmap.png}
    \caption{Stratified median clearance time by ship type and LOA bin for port and starboard transfer sides.}
    \label{fig:stratified_heatmap}
\end{figure}

\subsection{Operational use of percentile guidance}
The main operational value lies in percentile-based support rather than point prediction. Lower percentiles represent optimistic or minimally sufficient action durations. Upper percentiles provide conservative guidance for pilot MOB initiation. A practical rule is to use $P50$ as a minimum hold-time reference under routine high-readiness conditions, and $P90$ as a conservative hold-time envelope when uncertainty is elevated (reduced maneuverability, uncertain side geometry, or degraded bridge response). Bridge teams should treat these values as bounded initial-action windows, not deterministic forecasts.

\subsection{Sensitivity and model interpretation}
The one-at-a-time sensitivity analysis (Figure~\ref{fig:sensitivity_tornado}) indicates that the median clearance time is most responsive to hazard-radius inflation, reduced steering gain, and increased response delay. Two bars in the tornado chart perturb the same variable---crew response delay ($t_{\mathrm{bridge}}$)---by $+10$~s and $+30$~s respectively, bracketing the fixed 4.5~s baseline assumption and quantifying how operational readiness margins change under less favourable crew reaction conditions.

\begin{figure}[H]
    \centering
    \includegraphics[width=0.90\textwidth]{Figure_6_sensitivity_tornado.png}
    \caption{One-at-a-time sensitivity of the baseline median clearance time.}
    \label{fig:sensitivity_tornado}
\end{figure}

\section{Conclusion}
This study develops a practical framework for quantifying the first emergency action in pilot-transfer MOB response, before full recovery-turn reconstruction. By formulating immediate protection as a stern-clearance problem, the paper defines an initial-action envelope through two trigger variables: clearance time, $\tau_{\mathrm{clear}}$, and heading change at clearance, $\psi_{\mathrm{clear}}$. This variable pair is the primary conceptual contribution.

Within the current implementation, the executable results should be interpreted as a reproducible baseline rather than a fully calibrated estimate. The run uses operationally eligible accessible cases, applies an explicit hard-over fallback because maneuver-plan records do not align with the executable subset, adopts an IMO benchmark proxy for rudder travel, and activates wind-wave forcing for most cases while current remains unavailable. Even under these restrictions, the model yields a coherent percentile envelope and a strong side-dependent contrast.

The next stage should strengthen event-level data fusion for maneuver and metocean variables and extend validation against additional operational cases. The present paper therefore provides a formal definition and first quantified operational envelope for the pilot MOB initial-action problem.

\section*{CRediT authorship contribution statement}
To be completed.

\section*{Declaration of competing interest}
The authors declare no known competing interests.

\section*{Data availability}
Data and code are available from the authors upon reasonable request, subject to operational data-sharing constraints.

\section*{References}
\begin{thebibliography}{00}
\bibitem{ref_pilot_transfer_accidents}
Sakar, C., Turan, O., Incecik, A., 2024.
Dynamic analysis of pilot transfer accidents and associated safety recommendations.
\emph{Ocean Engineering} 293, 116718.
\url{https://doi.org/10.1016/j.oceaneng.2023.116718}

\bibitem{ref_ship_motion_review}
Zhang, L., Hong, Y., Wang, W., Wang, X., 2024.
Ship motion prediction methods and applications: A comprehensive review.
\emph{Journal of Marine Science and Engineering} 12(1), 107.
\url{https://doi.org/10.3390/jmse12010107}

\bibitem{ref_ship_motion_review2}
Cademartori, C., Moniaci, W., Viviani, M., 2023.
Review on ship motions and quiescent periods prediction methods.
\emph{Applied Ocean Research} 141, 103783.
\url{https://doi.org/10.1016/j.apor.2023.103783}

\bibitem{ref_ais_review}
Yang, Y., Li, Y., Liu, J., Zhang, Z., 2024.
A review of automatic identification system trajectory data in maritime applications.
\emph{Ocean Engineering} 306, 117995.
\url{https://doi.org/10.1016/j.oceaneng.2024.117995}

\bibitem{ref_aviation_alerts}
Fercho, M., Casner, S.M., Geven, R., 2024.
Pilots' response times and response strategies for resolving alerts in modern commercial aircraft.
\emph{Human Factors} 66(1), 162--178.
\url{https://doi.org/10.1177/00187208221115396}

\bibitem{ref_aviation_hud}
Zheng, Y., Wang, L., Hou, M., Pan, Z., 2024.
Pilots' reactions to collision warnings on head-up displays.
\emph{International Journal of Human-Computer Interaction} 40(9), 2595--2606.
\url{https://doi.org/10.1080/10447318.2023.2183943}

\bibitem{ref_aviation_envelope}
Zhang, Q., Wang, Y., de Visser, C.C., Chu, Q., 2021.
Database-driven safe flight-envelope protection for impaired aircraft.
\emph{Journal of Guidance, Control, and Dynamics} 44(3), 570--589.
\url{https://doi.org/10.2514/1.G005519}

\bibitem{ref_land_aeb}
Lai, F., Zhao, J., Zhang, S., 2024.
An active emergency braking control method based on time-to-collision and steering collision risk.
\emph{Machines} 12(1), 15.
\url{https://doi.org/10.3390/machines12010015}

\bibitem{ref_land_detection}
Fukui, M., Hayama, R., Takeda, K., Yamada, K., 2024.
The influence of pedestrian awareness of passenger cars on the driver's instinctive response timing in emergencies.
\emph{Transportation Research Part F: Traffic Psychology and Behaviour} 109, 515--531.
\url{https://doi.org/10.1016/j.trf.2024.10.010}

\bibitem{ref_camliyurt_pilot_fta2022}
Camliyurt, G., Choi, S.A., Kim, S.R., Guzel, A.T., et al., 2022.
Risk assessment for marine pilot occupational accidents using fault tree and event tree analysis.
\emph{Journal of Navigation and Port Research} (Korean), 2022.
\url{https://scholar.kyobobook.co.kr/article/detail/4010036927826}

\bibitem{ref_camliyurt_pilot_ml2023}
Camliyurt, G., Park, Y., Kim, D., Kang, W.S., et al., 2023.
Machine learning with multi-source data to predict and explain marine pilot occupational accidents.
\emph{Journal of Marine Science and Technology} 31(4), 2023.
\url{https://jmstt.ntou.edu.tw/journal/vol31/iss4/2/}

\bibitem{ref_yang_bn_maritime2014}
Li, K.X., Yin, J., Bang, H.S., Yang, Z., Wang, J., 2014.
Bayesian network with quantitative input for maritime risk analysis.
\emph{Maritime Policy \& Management} 41(2), 153--169.
\url{https://doi.org/10.1080/03088839.2013.863432}

\bibitem{ref_yang_psc_bn2018}
Yang, Z., Yang, Z., Yin, J., 2018.
Realising advanced risk-based port state control inspection using data-driven Bayesian networks.
\emph{Transportation Research Part A: Policy and Practice} 110, 38--56.
\url{https://doi.org/10.1016/j.tra.2018.02.003}

\bibitem{ref_yang_maritime_accident_bn2020}
Fan, S., Yang, Z., Blanco-Davis, E., et al., 2020.
Analysis of maritime transport accidents using Bayesian networks.
\emph{Proceedings of the Institution of Mechanical Engineers, Part O: Journal of Risk and Reliability} 234(5), 797--811.
\url{https://doi.org/10.1177/1748006X19900850}

\bibitem{ref_solas}
IMO, 1974 (as amended).
International Convention for the Safety of Life at Sea (SOLAS).
International Maritime Organization, London.

\bibitem{ref_iamsar}
IMO and ICAO, 2022.
International Aeronautical and Maritime Search and Rescue Manual (IAMSAR), Volume III: Mobile Facilities.
International Maritime Organization, London.

\bibitem{ref_fossen}
Fossen, T.I., 2021.
Handbook of Marine Craft Hydrodynamics and Motion Control, 2nd ed.
Wiley, Chichester.
\end{thebibliography}

\end{document}
